% Part: first-order-logic
% Chapter: semantics
% Section: intro-semantics

\documentclass[../../../include/open-logic-section]{subfiles}

\begin{document}

\olfileid{fol}{syn}{its}

\olsection{પરિચય}

અભિવ્યક્તિઓને અર્થ આપવાનું ક્ષેત્ર અર્થવિચાર છે. અર્થવિચારનો કેન્દ્રસ્થ
ખ્યાલ !!a{structure}માં સંતોષનો છે. !!^a{structure} ભાષાના પાયાના
ઘટકોને અર્થ આપે છે: !!a{domain} એટલે વસ્તુઓનો અરિક્ત ગણ. પરિમાણકોનું
અર્થઘટન આ વ્યાપકક્ષેત્ર પર વ્યાપ્ત થતા પરિમાણકો તરીકે થાય છે, !!{constant}sને
વ્યાપકક્ષેત્રના ઘટકો નિયુક્ત થાય છે, !!{function}sને !!{domain}માંથી પોતાનામાં
જતાં વિધેયો નિયુક્ત થાય છે અને !!{predicate}sને !!{domain} પરના સંબંધો
નિયુક્ત થાય છે. !!{domain} અને પાયાના સંકેતભંડોળને કરેલી નિયુક્તિઓ
મળીને !!a{structure} બને છે. !!^{variable}s, !!{formula}sમાં આવી શકે
છે અને અર્થવિચાર આપવા માટે તેમને પણ !!{domain}ના !!{element}s નિયુક્ત
કરવા પડે છે---આ ચલ-નિયુક્તિ છે. અંતે, સંતોષસંબંધ આ બધી બાબતોને એકત્ર
કરે છે. !!^a{formula}, !!a{structure}~$\Struct{M}$માં
!!a{variable} નિયુક્તિ~$s$ સાપેક્ષ સંતોષાઈ શકે છે; તેને
$\Sat{M}{!A}[s]$ લખીએ છીએ. આ સંબંધ પણ~$!A$ની રચના પરના અનુમાન વડે
વ્યાખ્યાયિત થાય છે: દાખલા તરીકે, તાર્કિક સંયોજકોની સત્યતા સારણીઓ
વાપરીને $(!A \land !B)$નો સંતોષ, $!A$ અને~$!B$ સંતોષાય છે (કે નથી)
તેના આધારે વ્યાખ્યાયિત થાય છે. પછી એવું નીકળે છે કે જો
!!{formula}~$!A$ !!a{sentence} હોય, એટલે તેમાં મુક્ત ચલો ન હોય, તો
!!{variable} નિયુક્તિનો ફેર પડતો નથી; તેથી આપણે !!{sentence}s કોઈ
!!{structure}sમાં સાદી રીતે સંતોષાય છે (કે નથી) એમ કહી શકીએ છીએ.

!!{sentence}s માટેના સંતોષસંબંધ~$\Sat{M}{!A}$ના આધારે પછી આપણે
પ્રામાણ્ય, ફલિતતા અને સંતોષ્યતાના પાયાના અર્થવિચારી ખ્યાલો
વ્યાખ્યાયિત કરી શકીએ છીએ. દરેક !!{structure}, !!^a{sentence}ને સંતોષે,
એટલે $\Entails !A$ હોય, તો તે પ્રમાણભૂત છે. !!{sentence}sના ગણમાંથી
તે ફલિત થાય, એટલે $\Gamma \Entails !A$ હોય, ત્યારે અને માત્ર ત્યારે
જ્યારે~$\Gamma$નાં બધાં !!{sentence}sને સંતોષતી દરેક !!{structure},
~$!A$ને પણ સંતોષે. અને કોઈ એક !!{structure}, !!{sentence}sના ગણમાંનાં
બધાં !!{sentence}sને એકસાથે સંતોષે, તો એ ગણ સંતોષ્ય છે. !!{formula}s
અનુમાનાત્મક રીતે વ્યાખ્યાયિત છે અને સંતોષ પણ !!{formula}sની રચના પરના
અનુમાન વડે વ્યાખ્યાયિત થાય છે; તેથી આપણા અર્થવિચારના ગુણધર્મો સાબિત
કરવા અને વ્યાખ્યાયિત અર્થવિચારી ખ્યાલોને પરસ્પર જોડવા આપણે અનુમાન
વાપરી શકીએ છીએ.

\end{document}
